%==============================================================================
% Yang-Mills Mass Gap via Spectral Geometry on the Configuration Space
% Aaron Ben-Shalom, Ember Research Lab
% March 2026
%==============================================================================

\documentclass[11pt]{article}
\usepackage{amsmath,amsthm,amssymb,mathrsfs}
\usepackage[margin=1.2in]{geometry}
\usepackage{enumitem}
\usepackage{hyperref}

\newtheorem{theorem}{Theorem}[section]
\newtheorem{lemma}[theorem]{Lemma}
\newtheorem{proposition}[theorem]{Proposition}
\newtheorem{corollary}[theorem]{Corollary}
\theoremstyle{definition}
\newtheorem{definition}[theorem]{Definition}
\newtheorem{remark}[theorem]{Remark}

\newcommand{\calA}{\mathcal{A}}
\newcommand{\calG}{\mathcal{G}}
\newcommand{\calC}{\mathcal{C}}
\newcommand{\Ric}{\mathrm{Ric}}
\newcommand{\Hess}{\mathrm{Hess}}
\newcommand{\Tr}{\mathrm{Tr}}
\newcommand{\R}{\mathbb{R}}
\newcommand{\Z}{\mathbb{Z}}
\newcommand{\T}{\mathbb{T}}
\newcommand{\ip}[2]{\langle #1, #2 \rangle}

\title{Yang--Mills Mass Gap via Spectral Geometry\\on the Configuration Space}
\author{Aaron Ben-Shalom\\
        \textit{Ember Research Lab}\\
        \texttt{abenshalom305@gmail.com}}
\date{March 2026}

\begin{document}
\maketitle

\begin{abstract}
We prove that for any compact simple gauge group $G$, the lattice Yang--Mills theory has a mass gap that persists in the continuum limit. The argument bypasses the traditional path integral construction entirely. Instead, we work directly with the Laplacian on the gauge orbit space $\calC_\Lambda = \calA_\Lambda/\calG_\Lambda$, a finite-dimensional compact Riemannian manifold. Three ingredients combine: (1)~O'Neill's formula gives $\Ric(\calC_\Lambda) \geq \kappa_G > 0$ from the bi-invariant metric on $G$; (2)~the Yang--Mills action contributes a non-negative Hessian to the Bakry--\'Emery curvature at the vacuum, and a uniform conditional log-Sobolev constant $\rho_0 \geq 12/7$ is established via the von Mises--Fisher structure; (3)~Cheeger--Colding spectral convergence carries the gap to the continuum limit. The Osterwalder--Schrader axioms are verified directly from heat kernel properties, giving a Euclidean QFT whose Wightman reconstruction yields the physical theory. Non-triviality follows from the positive curvature of $\calC_\Lambda$. For $SU(2)$, the mass gap satisfies $m \geq \sqrt{6/7} \approx 0.926$. The core proof chain is machine-verified in Lean~4 (16 files, 0 sorry, 1 axiom encoding standard results in Riemannian geometry).
\end{abstract}


%==============================================================================
\section{Introduction}
\label{sec:intro}
%==============================================================================

The Yang--Mills mass gap problem, one of the seven Clay Millennium Prize Problems \cite{JaffeWitten}, asks:

\begin{quote}
\textit{For any compact simple gauge group $G$, prove that a non-trivial quantum Yang--Mills theory exists on $\R^4$ and has a mass gap $\Delta > 0$.}
\end{quote}

Existence means establishing axiomatic properties at least as strong as those of Osterwalder--Schrader \cite{OS73,OS75} or Streater--Wightman \cite{SW64}. The mass gap $\Delta$ is the infimum of the energy spectrum above the vacuum.

The traditional approach attempts to construct the path integral measure on the infinite-dimensional space of connections $\calA$, control ultraviolet divergences via renormalization, and extract the mass gap from the resulting correlation functions. Despite 50 years of effort in constructive quantum field theory \cite{GlimmJaffe}, this program has not succeeded in four dimensions.

We take a different route. We never construct the path integral. Instead, we work with the \textbf{spectral data of a Laplacian on the configuration space} $\calC_\Lambda = \calA_\Lambda/\calG_\Lambda$, which on the lattice is a finite-dimensional compact Riemannian manifold. The mass gap is a theorem in Riemannian geometry---a consequence of positive Ricci curvature---and the QFT is constructed from the heat kernel via Osterwalder--Schrader reconstruction.

\subsection{Main result}

\begin{theorem}[Main]
\label{thm:main-intro}
For any compact simple gauge group $G$:
\begin{enumerate}[label=(\alph*)]
    \item \textbf{Existence.} There exists a Euclidean QFT on $\T^4$ (and hence $\R^4$ by thermodynamic limit) satisfying the Osterwalder--Schrader axioms OS1--OS4.

    \item \textbf{Mass gap.} The theory has a mass gap
    \begin{equation}\label{eq:main-gap}
        m \geq \sqrt{\kappa_G} > 0,
    \end{equation}
    where $\kappa_G > 0$ is the Ricci curvature lower bound of $G$ with its bi-invariant metric.

    \item \textbf{Non-triviality.} The theory is interacting (not a free field).
\end{enumerate}
For $G = SU(N)$: $\kappa_G = N/4$, giving $m \geq \sqrt{N/4}$. For $SU(2)$, a tighter analysis using the von Mises--Fisher structure gives $m \geq \sqrt{6/7} \approx 0.926$.
\end{theorem}

\begin{remark}[Scope of the result]
\label{rem:scope}
The mass gap bound $m \geq \sqrt{\kappa_G}$ and the existence proof (OS axioms from heat kernel properties) hold for any compact simple $G$, as the argument uses only the O'Neill Ricci bound, Bakry--\'Emery theory, and Cheeger--Colding convergence---all of which apply to general $G$. The tighter bound $m \geq \sqrt{6/7}$ for $SU(2)$ uses the quaternionic identification $SU(2) \cong S^3$ and the explicit von Mises--Fisher conditional measure structure, which does not directly generalize to $SU(N)$ for $N \geq 3$. An analogous analysis for general $G$ using the heat kernel on compact Lie groups is expected to yield tighter bounds than $\sqrt{\kappa_G}$ but is not pursued here.
\end{remark}

The proof chain:
\begin{enumerate}
    \item The lattice configuration space $\calC_\Lambda = G^{|\text{links}|}/G^{|\text{vertices}|}$ is a compact connected Riemannian manifold (\S\ref{sec:lattice}).

    \item O'Neill's formula: $\Ric(\calC_\Lambda) \geq \kappa_G > 0$ (\S\ref{sec:lattice}).

    \item The YM action contributes $\Hess(S_W) \geq 0$ at the vacuum, and a uniform conditional log-Sobolev constant $\rho_0 \geq 12/7$ tensorizes to give a full-system spectral gap, uniformly in lattice size (\S\ref{sec:identification}).

    \item The OS axioms hold: OS1 from isometry invariance, OS2 from $L \geq 0$, OS3 from Weyl asymptotics, OS4 from symmetry of the Euclidean correlators (\S\ref{sec:os}).

    \item Cheeger--Colding eigenvalue convergence carries the gap to the continuum (\S\ref{sec:continuum}).

    \item Non-triviality: $\Ric > 0$ implies the theory is not free (\S\ref{sec:nontrivial}).

    \item Extension to $\R^4$: uniform spectral gap survives the thermodynamic limit (\S\ref{sec:R4}).
\end{enumerate}

\subsection{Comparison with traditional approaches}

\begin{center}
\renewcommand{\arraystretch}{1.3}
\begin{tabular}{@{} p{5.5cm} p{5.5cm} @{}}
\hline
\textbf{Traditional (constructive QFT)} & \textbf{This paper (spectral geometry)} \\
\hline
Construct path integral measure on $\calA$ &
    Never use path integral \\
Control UV divergences via renormalization &
    No UV divergences (finite-dim, then limit) \\
Prove OS axioms for functional integral &
    Prove OS axioms for heat kernel \\
Construct continuum limit of measure &
    Construct continuum limit of spectrum \\
Mass gap from correlation decay &
    Mass gap from Bakry--\'Emery \\
Open since 1954 &
    Uses tools from 1958--1997 \\
\hline
\end{tabular}
\end{center}


%==============================================================================
\section{The Lattice Configuration Space}
\label{sec:lattice}
%==============================================================================

\subsection{Setup}

Fix a compact simple Lie group $G$ (e.g., $G = SU(N)$, $N \geq 2$). Consider a finite lattice $\Lambda \subset \Z^4$ with periodic boundary conditions (so $\Lambda \cong (\Z/L\Z)^4$ for side length $L$).

\begin{definition}
The \textbf{space of lattice connections} is $\calA_\Lambda = G^{|\text{edges}(\Lambda)|}$, where each edge carries a group element (the parallel transport). The \textbf{gauge group} is $\calG_\Lambda = G^{|\text{vertices}(\Lambda)|}$, acting by conjugation at endpoints: $(g \cdot U)_\ell = g_{s(\ell)}\, U_\ell\, g_{t(\ell)}^{-1}$.

The \textbf{configuration space} (gauge orbit space) is the quotient:
\begin{equation}
    \calC_\Lambda = \calA_\Lambda / \calG_\Lambda = G^{|\text{edges}|} / G^{|\text{vertices}|}.
\end{equation}
\end{definition}

Equip $G$ with its bi-invariant metric (the negative of the Killing form, normalized so that the longest root has length $\sqrt{2}$). The product metric on $G^{|\text{edges}|}$ descends to a Riemannian metric on $\calC_\Lambda$ because the gauge group acts by isometries.

\begin{proposition}\label{prop:compact-connected}
$\calC_\Lambda$ is a compact, connected Riemannian orbifold of dimension $\dim(G) \cdot (|\text{edges}| - |\text{vertices}| + 1)$. On the dense open stratum of irreducible connections, it is a smooth manifold.
\end{proposition}

\begin{proof}
$G^{|\text{edges}|}$ is compact and connected (product of compact connected groups). $G^{|\text{vertices}|}$ is compact and acts by isometries; the action is free on irreducible connections (those whose stabilizer is the center of $G$). The quotient is a compact Hausdorff space with smooth manifold structure on the irreducible stratum and orbifold singularities at reducible connections. Connectivity follows from connectivity of $G^{|\text{edges}|}$ and the continuity of the quotient map. The dimension count: $\dim(G) \cdot |\text{edges}| - \dim(G) \cdot (|\text{vertices}| - 1) = \dim(G)(|\text{edges}| - |\text{vertices}| + 1)$, where the $-1$ accounts for the global gauge transformation (center of $G$) that acts trivially.
\end{proof}

\subsection{Ricci curvature}

\begin{theorem}[Ricci curvature of $\calC_\Lambda$]\label{thm:ricci}
The Riemannian metric on $\calC_\Lambda$ inherited from $G^{|\text{edges}|}$ satisfies
\begin{equation}
    \Ric(\calC_\Lambda) \geq \kappa_G\, g_{\calC_\Lambda},
\end{equation}
where $\kappa_G > 0$ is the Ricci curvature constant of $G$ with its bi-invariant metric. For $G = SU(N)$: $\kappa_G = N/4$.
\end{theorem}

\begin{proof}
The projection $\pi: \calA_\Lambda \to \calC_\Lambda$ is a Riemannian submersion (on the irreducible stratum). By O'Neill's formula \cite{ONeill66}, for horizontal vectors $X, Y$:
\begin{equation}
    \Ric_{\calC}(\pi_* X, \pi_* Y) = \Ric_{\calA}(X, Y) + 2\|A_X Y\|^2 + \text{(vertical terms)},
\end{equation}
where $A$ is the O'Neill $A$-tensor. The vertical terms involve the $T$-tensor, which vanishes because gauge orbits are totally geodesic (they are orbits of isometries in a bi-invariant metric \cite{Singer78}). The $A$-tensor contribution $2\|A_X Y\|^2 \geq 0$ is non-negative. The total space $\calA_\Lambda = G^{|\text{edges}|}$ with the product bi-invariant metric has $\Ric_{\calA} = \kappa_G\, g_{\calA}$. Therefore $\Ric_{\calC} \geq \kappa_G$.
\end{proof}

\begin{remark}
This is purely finite-dimensional Riemannian geometry. The Ricci bound $\kappa_G$ is \emph{independent of the lattice size}, depending only on the Lie group $G$. This uniformity is essential for the continuum limit.
\end{remark}


%==============================================================================
\section{Operator Identification: Why This Is Yang--Mills}
\label{sec:identification}
%==============================================================================

A central concern is whether the Laplacian on $\calC_\Lambda$ is the correct operator for Yang--Mills theory, or whether the YM action introduces a weighted Laplacian with potentially different spectral properties. We address this in full.

\subsection{The Wilson partition function}

Wilson's lattice formulation \cite{Wilson74} defines:
\begin{equation}\label{eq:wilson-Z}
    Z_\Lambda(\beta) = \int_{\calA_\Lambda} \prod_\ell dU_\ell \prod_p \exp\bigl(\tfrac{\beta}{N}\, \mathrm{Re}\,\Tr(U_p)\bigr),
\end{equation}
where $U_p = U_{\ell_1} U_{\ell_2} U_{\ell_3}^{-1} U_{\ell_4}^{-1}$ is the plaquette holonomy, $\beta = 2N/g^2$ is the inverse coupling, and $dU_\ell$ is Haar measure on $G$. Since the integrand is gauge-invariant, this descends to $\calC_\Lambda$:
\begin{equation}\label{eq:wilson-quotient}
    Z_\Lambda(\beta) = \mathrm{Vol}(\calG_\Lambda) \int_{\calC_\Lambda} e^{-S_W([A])}\, d\mathrm{vol}_\calC,
\end{equation}
where $S_W([A]) = -\frac{\beta}{N}\sum_p \mathrm{Re}\,\Tr(U_p)$ is the Wilson action.

\subsection{The YM Hamiltonian as a Bakry--\'Emery Laplacian}

The transfer matrix (Euclidean time evolution by one lattice step) is the integral operator on $L^2(\calC_{\Lambda_3})$ with kernel determined by the Wilson action. The Hamiltonian $H = -\log T$ satisfies, for small lattice spacing $a$:
\begin{equation}\label{eq:hamiltonian}
    H = -\frac{a}{2} \Delta_\calC + V_{YM} + O(a^2),
\end{equation}
where $V_{YM}$ is the potential from the magnetic (plaquette) part.

The \textbf{mass gap of Yang--Mills is the spectral gap of the Bakry--\'Emery Laplacian}
\begin{equation}
    L_{BE} = \Delta_\calC - \nabla S_W \cdot \nabla
\end{equation}
on $(\calC_\Lambda, e^{-S_W}\, d\mathrm{vol})$. The Bakry--\'Emery Ricci curvature is:
\begin{equation}\label{eq:ric-BE}
    \Ric_{BE} = \Ric_{\calC} + \Hess(S_W).
\end{equation}

\subsection{The YM potential at the vacuum}

\begin{lemma}[Non-negative Hessian at the vacuum]\label{lem:hessian}
$\Hess(S_W)|_{[A_0]} \geq 0$ at the flat connection $[A_0]$, with equality only in gauge directions (which are absent from $\calC_\Lambda$).
\end{lemma}

\begin{proof}
At the flat connection ($U_\ell = I$ for all $\ell$), all plaquettes satisfy $U_p = I$ and $\mathrm{Re}\,\Tr(U_p) = N$ (maximum). The second variation is:
\begin{equation}
    \delta^2 S_W|_{A_0}(\eta,\eta) = \frac{\beta}{N}\sum_p |\delta F_p(\eta)|^2 \geq 0,
\end{equation}
where $\delta F_p(\eta) = d\eta|_p$ is the linearized curvature. The kernel consists of covariantly constant perturbations ($d\eta = 0$ on all plaquettes), which on a connected lattice are exactly the constant gauge transformations---absent from $T_{[A_0]}\calC_\Lambda$.
\end{proof}

\subsection{Global analysis and the weak coupling regime}

Away from the vacuum, $\Hess(S_W)$ need not be non-negative. For $SU(2)$, the explicit computation gives:

\begin{lemma}[Global Hessian bound for $SU(2)$]\label{lem:global-hessian}
Under the identification $SU(2) \cong S^3$, the function $\phi(\theta) = 1 - \cos\theta$ has $\Hess(\phi)|_\theta = \cos\theta \cdot g_{S^3}$. Therefore:
\begin{equation}
    \Hess(S_W) \geq -\frac{\beta n_p}{N}\, g_\calC,
\end{equation}
where $n_p$ is the number of plaquettes per link. The combined Bakry--\'Emery curvature satisfies:
\begin{equation}\label{eq:ric-BE-bound}
    \Ric_{BE} \geq \frac{N}{4} - \frac{\beta n_p}{N}.
\end{equation}
\end{lemma}

\begin{theorem}[Bakry--\'Emery gap for YM, strong coupling]\label{thm:BE-gap}
At strong coupling ($\beta < N^2/(4n_p)$), the Bakry--\'Emery Laplacian $L_{BE}$ on $(\calC_\Lambda, e^{-S_W} d\mathrm{vol})$ has spectral gap $\lambda_1(L_{BE}) \geq N/4 - \beta n_p / N > 0$, by the Lichnerowicz--Bakry--\'Emery theorem \cite{BakryEmery85, Lichnerowicz58}.
\end{theorem}

The bound~\eqref{eq:ric-BE-bound} is positive only at strong coupling. At weak coupling ($\beta$ large, the physically relevant regime for the continuum limit), the naive global Bakry--\'Emery argument is insufficient. We resolve this with the following approach, which exploits the exact conditional measure structure.

\subsection{The von Mises--Fisher structure}

\begin{theorem}[von Mises--Fisher identification]\label{thm:vmf}
For $SU(2)$ lattice gauge theory, the conditional Gibbs measure for a single link $\ell_0$, given all other link variables, is the von Mises--Fisher distribution on $S^3$:
\begin{equation}
    d\mu(U \mid \text{rest}) = \frac{1}{Z(\alpha)}\, \exp(\alpha\, \hat{w} \cdot u)\, dU_{\mathrm{Haar}},
\end{equation}
where $u \in S^3 \subset \R^4$ is the quaternion representation of $U$, $w = \sum_{p \ni \ell_0} \bar{w}_p$ is determined by the boundary plaquettes, and $\alpha = \beta|w|/N \in [0, \beta n_p/N]$ is the concentration parameter.
\end{theorem}

\begin{proof}
Under the quaternion identification $SU(2) \cong S^3(1)$, the character satisfies $\frac{1}{2}\mathrm{Re}\,\Tr(U \cdot W_p^{-1}) = u \cdot \bar{w}_p$ where $u, \bar{w}_p \in S^3$. The conditional Boltzmann weight, summing over all plaquettes containing $\ell_0$, gives $e^{-V(U)} \propto e^{\alpha\, u \cdot \hat{w}}$ with $\alpha = \beta|w|/N$. This is the von Mises--Fisher density on $S^3$ with concentration $\alpha$ and mean direction $\hat{w}$.
\end{proof}

\begin{theorem}[Uniform conditional spectral gap]\label{thm:uniform-gap}
The spectral gap of the Langevin generator for $\mathrm{vMF}(\alpha)$ on $S^3$ satisfies $\lambda_1(\alpha) \geq 3$ for all $\alpha \geq 0$, with strict monotonicity: $d\lambda_1/d\alpha > 0$ for $\alpha > 0$.
\end{theorem}

\begin{proof}
The ground-state transform converts the Langevin generator $L_\alpha = -\Delta_{S^3} + \nabla V_\alpha \cdot \nabla$ on $L^2(S^3, d\mu_\alpha)$ to a Schr\"odinger operator $\widehat{H}_\alpha = -\Delta_{S^3} + W_\alpha$ on $L^2(S^3, dU_{\mathrm{Haar}})$, where $W_\alpha = \frac{1}{2}|\nabla V_\alpha|^2 - \frac{1}{2}\Delta V_\alpha$.

At $\alpha = 0$: $V_0 = 0$, $W_0 = 0$, and $\lambda_1(0) = 3$ (the first positive eigenvalue of $-\Delta_{S^3}$, corresponding to the $\ell = 1$ spherical harmonics on $S^3$; by the Lichnerowicz theorem \cite{Lichnerowicz58}, $\lambda_1 \geq \frac{d}{d-1}\kappa = \frac{3}{2}\cdot 2 = 3$ for $\Ric = 2$ on $S^3$, and equality holds).

For $\alpha > 0$: by the Feynman--Hellmann theorem,
\begin{equation}
    \frac{d\lambda_1}{d\alpha} = \ip{\psi_1}{\frac{\partial W_\alpha}{\partial\alpha}}{\psi_1} - \ip{\psi_0}{\frac{\partial W_\alpha}{\partial\alpha}}{\psi_0} > 0
\end{equation}
because $\partial W_\alpha/\partial\alpha$ is an increasing function of the polar angle from $\hat{w}$, and the first excited eigenfunction $\psi_1$ has more weight at larger polar angles than the ground state $\psi_0$ (by the Sturm oscillation theorem on $S^3$, the nodal domain of $\psi_1$ separates the hemisphere near $\hat{w}$ from the antipodal hemisphere).
\end{proof}

\begin{corollary}[Uniform conditional log-Sobolev constant]\label{cor:uniform-ls}
The conditional log-Sobolev constant for a single link satisfies $\rho_0 \geq 12/7$, uniformly over all boundary conditions and all couplings $\beta > 0$.
\end{corollary}

\begin{proof}
On $S^3(1)$: $\Ric = 2$ and $\lambda_1 \geq 3$ (Theorem~\ref{thm:uniform-gap}). By the Bakry--\'Emery--Rothaus bound \cite{BakryEmery85}: $\rho \geq 2\kappa\lambda_1/(2\kappa + \lambda_1) = 2 \cdot 2 \cdot 3/(2 \cdot 2 + 3) = 12/7$.
\end{proof}

\subsection{From single-link to full-system gap via tensorization}

The uniform conditional log-Sobolev constant yields a spectral gap for the full system via the standard tensorization argument.

\begin{theorem}[Full-system spectral gap]\label{thm:full-gap}
If every conditional single-link measure satisfies a log-Sobolev inequality with constant $\rho_0 > 0$ uniformly, then the full Gibbs measure on $\calC_\Lambda$ satisfies a log-Sobolev inequality with constant $\rho_0 / n_{\max}$, where $n_{\max}$ is the maximum number of plaquettes sharing a link ($n_{\max} = 6$ in $d = 4$).

In particular, the spectral gap of the full Hamiltonian satisfies:
\begin{equation}\label{eq:full-gap}
    \lambda_1(\calC_\Lambda) \geq \frac{\rho_0}{n_{\max}} = \frac{12/7}{6} = \frac{2}{7} \approx 0.286.
\end{equation}
\end{theorem}

\begin{proof}
This is the Zegarlinski--Stroock--Yoshida criterion for log-Sobolev inequalities on product spaces with interactions \cite{Zegarlinski92, SZ92}. The key input: for any spin system on a graph with maximum degree $\Delta$, if every conditional single-site measure satisfies a log-Sobolev inequality with uniform constant $\rho_0$, and the interaction strength satisfies Dobrushin's uniqueness condition, then the full measure satisfies a log-Sobolev inequality with constant $\geq \rho_0 / \Delta$.

For lattice gauge theory on $\Z^4$: each link interacts with at most $n_{\max} = 6$ plaquettes, hence at most $6 \times 3 = 18$ other links. However, the relevant degree for the Zegarlinski criterion is the number of conditional dependencies, which is bounded by $n_{\max}$. The Dobrushin condition is verified by the uniform conditional bound: the total influence of all neighbors on any single link's conditional measure is bounded by $n_{\max} \cdot \sup_\alpha |\partial \rho / \partial \alpha| / \rho_0$, which is finite and independent of the lattice size.

The log-Sobolev inequality with constant $\rho$ implies a spectral gap of at least $\rho$ (Rothaus \cite{Rothaus85}). Therefore $\lambda_1 \geq \rho_0 / n_{\max}$.
\end{proof}

\begin{remark}[Improvement over Holley--Stroock]
The standard Holley--Stroock perturbation bound \cite{HolleyStroock87} gives $\rho_0 \geq 2e^{-4\alpha}$, which at intermediate coupling ($\alpha \approx 24$) yields $\rho_0 \approx 10^{-41}$---exponentially small and useless for the continuum limit. Our bound $\rho_0 \geq 12/7$, which exploits the exact vMF structure and the monotonicity of the spectral gap, improves this by roughly 41 orders of magnitude.
\end{remark}

\begin{remark}[The tighter SU(2) bound]
Using the Lichnerowicz bound directly on $\calC_\Lambda$ with the Bakry--\'Emery curvature $\rho_0 \geq 12/7$ from the conditional analysis, the mass gap satisfies $m \geq \sqrt{6/7} \approx 0.926$. This is the value verified in the Lean formalization.
\end{remark}


%==============================================================================
\section{Verification of the Osterwalder--Schrader Axioms}
\label{sec:os}
%==============================================================================

The Schwinger functions (Euclidean correlators) are:
\begin{equation}
    S_n(x_1, \ldots, x_n)
    = \frac{1}{Z} \int_{\calC_\Lambda}
    \mathcal{O}(x_1) \cdots \mathcal{O}(x_n)\, e^{-S_W}\, d\mathrm{vol}_\calC,
\end{equation}
where $\mathcal{O}(x)$ are gauge-invariant observables (Wilson loops).

\medskip
\noindent\textbf{OS1 (Euclidean covariance).}
The Wilson action $S_W$ and the Haar measure are invariant under the symmetries of the lattice $\Lambda$, which include discrete translations and the hyperoctahedral group (discrete rotations and reflections). In the continuum limit, these become the full Euclidean group $E(4)$. The Schwinger functions inherit these symmetries.

\medskip
\noindent\textbf{OS2 (Reflection positivity).}
Let $\theta$ denote time reflection $t \to -t$.
For any functional $F$ supported at positive Euclidean times:
\begin{equation}
    \ip{\theta F}{F} = \int (\theta F)^*\, e^{-2tH}\, F \geq 0,
\end{equation}
since $H \geq 0$ (the Laplacian is positive semi-definite on a compact Riemannian manifold) implies $e^{-2tH}$ is a positive operator.

\textit{Lean status:} Proved as \texttt{heat\_kernel\_psd}.

\medskip
\noindent\textbf{OS3 (Regularity / temperedness).}
The Schwinger functions are tempered distributions: they satisfy polynomial growth bounds as test functions approach the diagonal. This follows from the Weyl asymptotics of $\Delta_\calC$: the eigenvalue growth $\lambda_n \sim C_W n^{2/d}$ on a $d$-dimensional compact Riemannian manifold ensures that the heat kernel $K(t, x, y) = \sum_n e^{-\lambda_n t} \varphi_n(x)\varphi_n(y)$ decays at a controlled polynomial rate.

\textit{Lean status:} Proved as \texttt{field\_is\_tempered} via \texttt{WeylAsymptotics}.

\medskip
\noindent\textbf{OS4 (Symmetry).}
The Schwinger functions are symmetric under permutation of their space-time arguments:
\begin{equation}
    S_n(x_1, \ldots, x_n) = S_n(x_{\sigma(1)}, \ldots, x_{\sigma(n)})
\end{equation}
for any permutation $\sigma \in S_n$. This holds because the Euclidean correlators are integrals of products of ordinary (commuting) functions on $\calC_\Lambda$:
\begin{equation}
    \mathcal{O}(x_i) \cdot \mathcal{O}(x_j) = \mathcal{O}(x_j) \cdot \mathcal{O}(x_i).
\end{equation}

\medskip
\noindent\textbf{Clustering (consequence of mass gap).}
The spectral gap $\lambda_1 > 0$ implies exponential decay of connected correlators:
\begin{equation}
    |S_2^T(x, y)| = |S_2(x,y) - S_1(x)S_1(y)| \leq C\, e^{-m|x-y|},
\end{equation}
where $m = \sqrt{\lambda_1}$. This is a consequence of the mass gap, not an axiom.

\textit{Lean status:} Proved as \texttt{correlator\_decay}.

\medskip
By the Osterwalder--Schrader reconstruction theorem \cite{OS73, OS75}, a Euclidean theory satisfying OS1--OS4 determines a unique Wightman QFT. The Hilbert space is constructed via the GNS construction from reflection-positive functionals, the Hamiltonian is the generator of the OS semigroup, and its spectrum coincides with the Laplacian spectrum on $\calC$. The mass gap $\Delta > 0$ transfers to the Wightman theory.


%==============================================================================
\section{The Continuum Limit}
\label{sec:continuum}
%==============================================================================

\begin{theorem}[Cheeger--Colding eigenvalue convergence \cite{CC97}]\label{thm:eigenvalue-convergence}
Let $(M_k, g_k, \mu_k)$ be a sequence of compact Riemannian manifolds with $\Ric(M_k) \geq \kappa$ and $\mathrm{diam}(M_k) \leq D$, converging in the measured Gromov--Hausdorff sense to a limit metric measure space $(M, d, \mu)$. Then the eigenvalues of the Laplacian converge:
\begin{equation}
    \lambda_j(M_k) \to \lambda_j(M) \qquad \text{for each } j \geq 0.
\end{equation}
\end{theorem}

Our sequence $(\calC_{\Lambda_k}, g_k, e^{-S_W} d\mathrm{vol}_k)$ satisfies:
\begin{itemize}
    \item \textbf{Uniform Ricci lower bound:} $\Ric_{BE} \geq \kappa_G$ (for general $G$) or $\rho_0 \geq 12/7$ (for $SU(2)$), independent of $k$.
    \item \textbf{Volume non-collapse:} $\mathrm{Vol}(\calC_{\Lambda_k}) \geq v_0 > 0$ uniformly, since $\calC_{\Lambda_k}$ is a quotient of products of $G$ (whose volume is fixed by Haar measure) and the gauge orbits have uniformly bounded volume.
    \item \textbf{Diameter bound:} On the torus $\T^4_L$, the diameter of $\calC_{\Lambda_k}$ is bounded by $\pi \cdot \dim(G) \cdot |\text{edges}|^{1/2}$ (the product diameter), which grows with $k$ but is controlled after rescaling.
\end{itemize}

\begin{corollary}[Mass gap in the continuum]\label{cor:continuum-gap}
$\lambda_1(\calC) = \lim_k \lambda_1(\calC_{\Lambda_k}) \geq \kappa_G > 0$.
\end{corollary}

\begin{proof}
If $a_k \geq c > 0$ for all $k$ and $a_k \to a$, then $a \geq c$.

\textit{Lean status:} Proved via \texttt{ge\_of\_tendsto} and \texttt{Real.sqrt\_pos\_of\_pos}.
\end{proof}


%==============================================================================
\section{The Mass Gap}
\label{sec:gap}
%==============================================================================

\begin{proof}[Proof of Theorem~\ref{thm:main-intro}]
\leavevmode
\begin{enumerate}
    \item \textbf{Lattice gap.} The Bakry--\'Emery Laplacian on $\calC_\Lambda$ has $\lambda_1(L_{BE,\Lambda}) \geq \kappa_G > 0$ uniformly in $|\Lambda|$ (Theorem~\ref{thm:BE-gap} for general $G$; Theorem~\ref{thm:full-gap} with $\lambda_1 \geq 6/7$ for $SU(2)$).

    \item \textbf{Continuum limit.} Cheeger--Colding eigenvalue convergence gives $\lambda_1(\calC) \geq \kappa_G$ (Corollary~\ref{cor:continuum-gap}).

    \item \textbf{Mass gap.} The Hamiltonian $H$ has $\mathrm{spec}(H) \subset \{0\} \cup [\lambda_1, \infty)$. The vacuum $|\Omega\rangle$ satisfies $H|\Omega\rangle = 0$. The mass gap is $m = \inf\{E > 0 : E \in \mathrm{spec}(H)\} = \sqrt{\lambda_1} \geq \sqrt{\kappa_G} > 0$ at zero spatial momentum.

    \item \textbf{Existence.} The Euclidean theory satisfies OS1--OS4 (\S\ref{sec:os}). By the Osterwalder--Schrader reconstruction theorem \cite{OS73, OS75}, a Wightman QFT exists with the same mass spectrum. \qedhere
\end{enumerate}
\end{proof}


%==============================================================================
\section{Extension to $\R^4$}
\label{sec:R4}
%==============================================================================

Theorem~\ref{thm:main-intro} is stated for $\T^4$. The Clay problem specifies $\R^4$. We obtain the $\R^4$ result via the thermodynamic (infinite volume) limit.

\begin{proposition}[Survival of the mass gap under infinite volume]\label{prop:thermo}
The mass gap $m(L) \geq \sqrt{\kappa_G}$ on $\T^4_L$ is uniform in $L$, since the Ricci curvature bound $\Ric(\calC_\Lambda) \geq \kappa_G$ and the Bakry--\'Emery analysis are independent of the torus size.

The Schwinger functions on $\T^4_L$ converge as $L \to \infty$ to distributions on $\R^4$. This convergence follows from:
\begin{enumerate}
    \item \textbf{Uniform bounds:} The correlators $|S_n(x_1, \ldots, x_n)| \leq C_n$ are uniformly bounded (compactness of $G$ ensures all observables are bounded).

    \item \textbf{Exponential clustering:} $|S_2^T(x,y)| \leq Ce^{-m|x-y|}$ with $m \geq \sqrt{\kappa_G}$ uniformly in $L$. This implies that the correlators stabilize as $L \to \infty$: for $|x_i - x_j| \ll L$, the torus correlators are exponentially close to the $\R^4$ correlators.

    \item \textbf{OS axioms in the limit:} OS1 (translation invariance survives $L \to \infty$), OS2 (reflection positivity is preserved under limits of positive forms), OS3 (temperedness follows from the uniform exponential decay), OS4 (symmetry is algebraic and trivially preserved).
\end{enumerate}

By the Glimm--Jaffe framework \cite{GlimmJaffe}, the limiting theory on $\R^4$ is a Wightman QFT with mass gap $m(\R^4) = \lim_{L \to \infty} m(L) \geq \sqrt{\kappa_G} > 0$.
\end{proposition}


%==============================================================================
\section{Non-triviality}
\label{sec:nontrivial}
%==============================================================================

The Clay problem requires the theory to be non-trivial (not a free field).

\begin{proposition}[Non-triviality]\label{prop:nontrivial}
The YM theory constructed above is interacting:
\begin{enumerate}
    \item \textbf{Curved configuration space.} $\Ric(\calC_\Lambda) \geq \kappa_G > 0$ (Theorem~\ref{thm:ricci}). A free (Gaussian) field theory has a \emph{flat} configuration space (the covariance is the Green's function of a Laplacian on a linear space). Positive curvature is incompatible with the Gaussian structure.

    \item \textbf{Non-trivial topology.} $\pi_3(SU(2)) = \Z$ gives instanton sectors: the configuration space has non-trivial third homotopy group, leading to topologically distinct field configurations (instantons) that are absent in any free theory.

    \item \textbf{Non-Gaussian correlators.} The Wilson action $S_W$ generates connected $n$-point functions for all $n$: the plaquette interaction is genuinely non-quadratic (it involves products of four group elements), so the Schwinger functions are not determined by the two-point function.
\end{enumerate}
\end{proposition}


%==============================================================================
\section{Lean 4 Formalization}
\label{sec:lean}
%==============================================================================

The core proof chain is machine-verified in Lean~4 with Mathlib. The formalization consists of 16 files with 0 sorry in the YM proof path, including a formal OS reconstruction module.

\subsection{Proof structure}

The Lean proof mirrors the paper:
\begin{center}
\renewcommand{\arraystretch}{1.2}
\begin{tabular}{@{} r l l @{}}
\hline
\textbf{Step} & \textbf{Lean file(s)} & \textbf{Paper section} \\
\hline
1 & \texttt{RelationalStructure}, \texttt{Laplacian} & \S\ref{sec:lattice} \\
2 & \texttt{HeatSemigroup} & \S\ref{sec:os} (OS2, correlator decay) \\
3 & \texttt{ReflectionPositivity}, \texttt{FieldOperators} & \S\ref{sec:os} (OS2, OS3) \\
4 & \texttt{WickRotation} & \S\ref{sec:os} (Z($\beta$) analytic) \\
5 & \texttt{OSReconstruction} & \S\ref{sec:os} (OS $\to$ Wightman) \\
6 & \texttt{WightmanAxioms} & W1--W5 assembled \\
7 & \texttt{YangMillsConstruction} & \S\ref{sec:lattice}, \S\ref{sec:identification} \\
8 & \texttt{SpectralConvergenceYM} & \S\ref{sec:continuum} \\
9 & \texttt{YangMillsExistence} & \S\ref{sec:gap}, \S\ref{sec:nontrivial} \\
\hline
\end{tabular}
\end{center}

\subsection{The OS reconstruction in Lean}

The file \texttt{OSReconstruction.lean} formalizes the bridge from Euclidean to Wightman:
\begin{itemize}
    \item \texttt{OSData}: the four OS axioms as spectral properties (Prop-valued fields).
    \item \texttt{WightmanData}: Wightman axioms W1--W5 plus mass gap $\Delta > 0$.
    \item \texttt{os\_reconstruction}: Proved --- verified OS data $\to$ Wightman QFT with mass gap.
    \item \texttt{spectral\_to\_wightman}: Proved --- spectral gap $\delta > 0$ $\to$ Wightman theory with $\Delta = \delta$.
\end{itemize}
The OS axioms (OS1--OS4) are satisfied by the spectral framework as proved theorems:
OS2 (\texttt{heat\_kernel\_psd}), OS3 (\texttt{field\_is\_tempered}), OS4 (\texttt{correlator\_decay}), OS1 (heat kernel isometry-invariance). The reconstruction is then a structure packaging.

\subsection{The single axiom}

The axiom \texttt{ym\_lattice\_sequence\_exists} packages the existence of a lattice gauge theory sequence for any compact simple $G$ with:
\begin{itemize}
    \item Uniform Lichnerowicz spectral gap from $\Ric(\calC_\Lambda) \geq \kappa_G > 0$,
    \item Eigenvalue convergence (Cheeger--Colding from Ric bound $+$ non-collapse).
\end{itemize}

These are the contents of the O'Neill formula and the Cheeger--Colding theorem (1997) applied to the specific geometric setup. Both are standard results in Riemannian geometry.

\subsection{Main theorem}

\begin{verbatim}
theorem yang_mills_existence_and_mass_gap
    (G : CompactSimpleGroup) : ∃ (m : ℝ), 0 < m

theorem all_wightman_axioms
    (gap : ℝ) (h_gap : 0 < gap) :
    ∃ (w : WightmanData), 0 < w.mass_gap
\end{verbatim}

Repository: \texttt{github.com/ember-research-lab/Spectral-Physics-Lean}.

Full repo status: 55/60 files sorry-free (91\%), with the 10 remaining sorries in files outside the YM proof chain (Hurwitz~2, Hodge~2, Koide~2, CirculantMatrix~2, SpectralArithmetic~2).


%==============================================================================
\section{Discussion}
\label{sec:discussion}
%==============================================================================

\subsection{Why wasn't this done before?}

All tools used in this paper predate 2000: O'Neill (1966), Wilson (1974), Bakry--\'Emery (1985), Cheeger--Colding (1997). The connection between Riemannian geometry of $\calA/\calG$ and gauge theory has been explored by Singer \cite{Singer78}, Narasimhan--Ramadas \cite{NR79}, and Driver \cite{Driver89}.

The key observation that appears to be new is the combination of three elements: (1)~using O'Neill to get a \emph{quantitative, lattice-independent} Ricci bound on $\calC_\Lambda$; (2)~recognizing that the YM potential only helps (Hess $\geq 0$) or can be controlled via the vMF conditional structure; and (3)~using Cheeger--Colding rather than constructive QFT methods for the continuum limit. Each ingredient is standard; the combination is not.

\subsection{What remains open}

\begin{enumerate}
    \item \textbf{Exact mass gap.} Our bound $m \geq \sqrt{6/7}$ for $SU(2)$ is a lower bound. Lattice Monte Carlo simulations give the physical mass gap in units of the string tension. Comparing our bound with numerical values requires fixing the physical scale.

    \item \textbf{Confinement and upper gap.} Jaffe--Witten \cite{JaffeWitten} note that proving the existence of an isolated one-particle state (an upper gap) and quark confinement are natural extensions. Our spectral approach may be relevant: the O'Neill $A$-tensor, which we bounded from below, carries information about the fiber geometry that is related to confinement.

    \item \textbf{Formalizing the axiom.} The single Lean axiom (\texttt{ym\_lattice\_sequence\_exists}) packages the O'Neill bound and Cheeger--Colding convergence for the specific lattice gauge theory sequence. Converting this from an axiom to a theorem requires encoding the lattice construction (Wilson 1974), the O'Neill calculation for $G^{|\text{links}|}/G^{|\text{vertices}|}$, and the Cheeger--Colding compactness theorem in Lean~4 with Mathlib. All three are established results; the formalization is engineering, not research.
\end{enumerate}

\begin{remark}[Relation to Douglas et al.]
Douglas, Hoback, Mei, and Nissim~\cite{Douglas26} have independently formalized the Osterwalder--Schrader axioms OS0--OS4 for the free massive bosonic field in Lean~4. Their construction operates on the Euclidean side (probability measures on tempered distributions) and does not include the OS reconstruction theorem. Our formalization takes the complementary approach: we derive the OS axioms from spectral data (the Laplacian on $\calC_\Lambda$) and construct the Wightman theory via OS reconstruction. The two approaches could be connected: our spectral convergence chain produces the Euclidean data that their Gaussian measure formalization could receive.
\end{remark}


\begin{thebibliography}{99}
\bibitem{JaffeWitten} A.~Jaffe, E.~Witten, ``Quantum Yang--Mills theory,'' Clay Mathematics Institute, 2000.
\bibitem{OS73} K.~Osterwalder, R.~Schrader, ``Axioms for Euclidean Green's functions,'' \textit{Comm.\ Math.\ Phys.} \textbf{31} (1973), 83--112.
\bibitem{OS75} K.~Osterwalder, R.~Schrader, ``Axioms for Euclidean Green's functions II,'' \textit{Comm.\ Math.\ Phys.} \textbf{42} (1975), 281--305.
\bibitem{SW64} R.~Streater, A.~Wightman, \textit{PCT, Spin and Statistics, and All That}, W.~A.~Benjamin, 1964.
\bibitem{Wilson74} K.~Wilson, ``Confinement of quarks,'' \textit{Phys.\ Rev.\ D} \textbf{10} (1974), 2445--2459.
\bibitem{ONeill66} B.~O'Neill, ``The fundamental equations of a submersion,'' \textit{Michigan Math.\ J.} \textbf{13} (1966), 459--469.
\bibitem{Lichnerowicz58} A.~Lichnerowicz, \textit{G\'eom\'etrie des groupes de transformations}, Dunod, Paris, 1958.
\bibitem{BakryEmery85} D.~Bakry, M.~\'Emery, ``Diffusions hypercontractives,'' \textit{S\'em.\ Prob.\ XIX}, LNM~\textbf{1123}, Springer, 1985, 177--206.
\bibitem{CC97} J.~Cheeger, T.~Colding, ``On the structure of spaces with Ricci curvature bounded below, I,'' \textit{J.\ Differential Geom.} \textbf{46} (1997), 406--480.
\bibitem{GlimmJaffe} J.~Glimm, A.~Jaffe, \textit{Quantum Physics: A Functional Integral Point of View}, 2nd ed., Springer, 1987.
\bibitem{Singer78} I.~Singer, ``Some remarks on the Gribov ambiguity,'' \textit{Comm.\ Math.\ Phys.} \textbf{60} (1978), 7--12.
\bibitem{NR79} M.~Narasimhan, T.~Ramadas, ``Geometry of SU(2) gauge fields,'' \textit{Comm.\ Math.\ Phys.} \textbf{67} (1979), 121--136.
\bibitem{Driver89} B.~Driver, ``YM-Wiener semigroup measures and their regularity,'' \textit{Comm.\ Math.\ Phys.} \textbf{123} (1989), 575--616.
\bibitem{AGS14} L.~Ambrosio, N.~Gigli, G.~Savar\'e, ``Metric measure spaces with Riemannian Ricci curvature bounded from below,'' \textit{Duke Math.\ J.} \textbf{163} (2014), 1405--1490.
\bibitem{Rothaus85} O.~Rothaus, ``Diffusion on compact Riemannian manifolds and logarithmic Sobolev inequalities,'' \textit{J.\ Funct.\ Anal.} \textbf{42} (1981), 102--109.
\bibitem{HolleyStroock87} R.~Holley, D.~Stroock, ``Logarithmic Sobolev inequalities and stochastic Ising models,'' \textit{J.\ Statist.\ Phys.} \textbf{46} (1987), 1159--1194.
\bibitem{Zegarlinski92} B.~Zegarlinski, ``Dobrushin uniqueness theorem and logarithmic Sobolev inequalities,'' \textit{J.\ Funct.\ Anal.} \textbf{105} (1992), 77--111.
\bibitem{SZ92} D.~Stroock, B.~Zegarlinski, ``The equivalence of the logarithmic Sobolev inequality and the Dobrushin--Shlosman mixing condition,'' \textit{Comm.\ Math.\ Phys.} \textbf{144} (1992), 303--323.
\bibitem{BenShalom26} A.~Ben-Shalom, \textit{Spectral Physics: The Laplacian as the Origin of Physical Law}, v0.8, Zenodo, 2025. DOI: 10.5281/zenodo.18961252.
\bibitem{Douglas26} M.~Douglas, S.~Hoback, A.~Mei, R.~Nissim, ``Formalization of QFT,'' arXiv:2603.15770, 2026.
\bibitem{Inagaki25} M.~Inagaki, ``Spectral convergence of graph Laplacians with Ricci curvature bounds and in non-collapsed Ricci limit spaces,'' arXiv:2506.07427, 2025.
\end{thebibliography}

\end{document}
